\section{总结与延伸：四个不等式，以及它们的适用条件}

讲座最后用四个简洁比较压缩主线：reasoning 优于 no reasoning，RL finetuning 优于 SFT，聚合多个答案优于单答案，retrieval + reasoning 优于单独 reasoning。它们应读成\textbf{在讲座所讨论任务和预算内的设计方向}，而不是无条件数学定理。每个“更好”都交换了成本、基础设施与新的失败模式。

\subsection{课堂四条结论怎样落到系统栈}

slide 47 是一张非常紧凑的总结页，本节把它展开成可实现的系统栈。读图时先看四个“大于号”分别修改推理长度、参数分布、采样统计和外部信息，再检查每一层新增了什么成本与验证责任。四种干预可以叠加在同一个系统中，而不是四选一。

\lecturefigure{slide-47.jpg}{课堂总结：reasoning、RL finetuning、aggregation 与 retrieval}{47}

\begin{center}
\begin{tabularx}{0.97\textwidth}{p{0.18\textwidth}p{0.23\textwidth}X X}
\toprule
干预层 & 改变什么 & 主要收益 & 新风险 \\
\midrule
Reasoning tokens & 单次生成的串行计算 & 分解、临时状态、可检查步骤 & 延迟、冗余、错误累积 \\
RL finetuning & 轨迹与答案的基础分布 & 成功路径更常出现 & verifier 偏差、reward hacking \\
Aggregation & 多次采样后的答案估计 & 降低单路径偶然性 & token 成本、相关错误、解析 \\
Retrieval & 模型可见的外部证据 & 时效知识、案例、原则 & 噪声、过期、来源与注入风险 \\
\bottomrule
\end{tabularx}
\end{center}

\begin{importantbox}{一张式子概括整套系统}
一个实用 reasoning system 可以写成：先检索 $z$，再从策略采样多条轨迹 $r^{(1:K)}$ 与答案 $a^{(1:K)}$，用 verifier 和 aggregation 选择最终输出。能力来自组件组合；可靠性取决于最弱接口是否可观测、可验证、可回放。
\end{importantbox}

\subsection{下一次突破：超越唯一可验证答案}

slide 48 把开放问题指向两处：解决没有唯一可验证答案的任务，以及建设真实应用而不是继续饱和 benchmark。前者要求从 binary correctness 走向多目标评价、过程证据、长期反馈和不确定性表达；后者要求端到端记录用户价值、失败严重性、成本、延迟、权限与可恢复性。

\lecturefigure{slide-48.jpg}{下一阶段：非唯一可验证任务与真实应用}{48}

\teachervoice{00:56:31--00:57:30，老师说 benchmark 很快会饱和，更想看到真实应用和超越唯一可验证答案的任务。课堂提示：这不是否定 benchmark，而是要求它们成为诊断工具，而非最终产品价值的替身。}

开放任务可以采用多层验收，而不是强迫一个 reward 承担全部价值：
\begin{enumerate}
  \item \textbf{硬约束}：事实来源、权限、格式、安全与法律边界；
  \item \textbf{任务指标}：完整性、可执行性、用户目标与成本；
  \item \textbf{过程证据}：检索来源、工具 trace、反例与不确定性；
  \item \textbf{长期结果}：真实采用、纠错率、回滚与负面外部性。
\end{enumerate}

\subsection{Q\&A：课堂结论不能删掉的六个限定}

讲座后的九分钟问答补上了 slides 中最重要的不确定性。第一，answer-token log-prob 在讲者经验中可作为 hallucination 信号，但没有被表述为完美 detector。第二，search 可以作为模型调用的工具；“研究 reasoning”时排除外部 search，不等于产品不该用 search。第三，训练如何重塑完整序列分布仍缺少令人满意的解释，不能把 top-token 直觉扩展成完整理论。

\teachervoice{00:57:30--01:01:00，老师回答 confidence 与 search：log-prob 在经验上有用，但 search 是可集成工具；他关心的是不依赖显式穷举时模型本身能学到什么。讲义提醒：研究隔离变量与产品追求最优性能是两个目标。}

第四，reasoning 与 answer 的边界需要 parser；若最终答案是一段程序，抽取和验证会明显更难。第五，低置信度候选有时可能正确，self-consistency 也会失败。第六，面对“AGI 两到五年、孩子该学什么”的问题，老师拒绝接受时间表前提，并强调当前方法仍有许多限制，真正缺少的是 killer applications；他明确认可编程助手价值，但没有把聊天产品等同于已到来的 AGI。

\teachervoice{01:01:00--01:06:04，老师承认 distribution reshaping 尚不清楚、程序答案需要仔细解析、self-consistency 不完美，并对短期 AGI 时间表保持怀疑。最后的实践标准是：模型是否形成真正有用的应用，而不是讨论热度。}

\begin{warningbox}{2025 课堂证据边界}
本讲中的 Gemini 2.0 thinking、o1 aggregation、Deep Research UI、GSM8K 数字与产品判断都应按 2025 年 4 月课堂时间点理解。讲义保留机制与工程启示，不把当时截图和经验陈述改写成 2026 年当前产品规格或普遍保证。
\end{warningbox}

\subsection{启动 reasoning 项目前的十二问}

\small
\begin{multicols}{2}
\begin{enumerate}
  \item “Reasoning” 在本任务中具体指哪些中间状态？
  \item 直接回答失败，是能力缺失还是 decoding 排序失败？
  \item 答案边界能否稳定解析，等价答案怎样规范化？
  \item 人工轨迹是否覆盖模型真正会访问的状态？
  \item Reward 与真实质量之间有什么代理差距？
  \item Verifier 的 false positive 和 false negative 多大？
  \item 扩展的是轨迹长度、样本数、搜索树还是工具调用？
  \item 多样本是否足够独立，还是共享同一系统偏差？
  \item Consistency 是否在目标分布上经过校准？
  \item Retrieval 的来源、时间与冲突怎样保存？
  \item 开放任务如何组合硬约束、偏好与长期反馈？
  \item 上线后如何回放 trace、发现回归并停止或回滚？
\end{enumerate}
\end{multicols}
\normalsize

\subsection{拓展阅读}

\scriptsize
\begin{multicols}{2}
\noindent\textbf{串行计算：}Li et al., \emph{Chain of Thought Empowers Transformers to Solve Inherently Serial Problems}. \href{https://arxiv.org/abs/2402.12875}{2402.12875}.\par
\noindent\textbf{无提示解码：}Wang and Zhou, \emph{Chain-of-Thought Reasoning Without Prompting}. \href{https://arxiv.org/abs/2402.10200}{2402.10200}.\par
\noindent\textbf{答案聚合：}Wang et al., \emph{Self-Consistency Improves Chain of Thought Reasoning in Language Models}. \href{https://arxiv.org/abs/2203.11171}{2203.11171}.\par
\noindent\textbf{自由文本聚合：}Chen et al., \emph{Universal Self-Consistency for Large Language Model Generation}. \href{https://arxiv.org/abs/2311.17311}{2311.17311}.\par
\end{multicols}
\normalsize

\end{document}
